\documentclass[lettersize,journal]{IEEEtran}
\usepackage{amsmath,amsfonts,amssymb}
\usepackage{algorithmic}
\usepackage{algorithm}
\usepackage{array}
\usepackage[caption=false,font=normalsize,labelfont=sf,textfont=sf]{subfig}
\usepackage{textcomp}
\usepackage{stfloats}
\usepackage{url}
\usepackage{verbatim}
\usepackage{graphicx}
\usepackage{color}
% \usepackage{appendix}
\usepackage{cite}
\newcommand{\parab}[1]{\noindent\textbf{#1} }
\newenvironment{proof}{{\noindent\it Proof:} }{\hfill $\square$\par}
\hyphenation{op-tical net-works semi-conduc-tor IEEE-Xplore} 

\begin{document}
\subsection{Sharper bounds for the finite profiles}
\label{app:tight-profile-bounds}
The norm choices in Eqs.~(1) and~(4) of the supplementary
proofs for Sections~V and~VI are sufficient choices for the
general parameter family. For the finite profiles we use the following sharper
bounds for the same distributions and norm tests. They replace the
conservative numerical sufficient conditions, without changing KeyGen or
the security games.

Write $\rho_t(c)=\sum_{z\in\mathbb Z}\exp(-(z-c)^2/(2t^2))$ and
$\rho_t=\rho_t(0)$. Poisson summation gives
$\rho_t(c)\le\rho_t$ and, for $t'\ge t$,
$\rho_{t'}/\rho_t\le t'/t$. The first inequality follows by bounding
each cosine in the dual series by one. For the second, the dual series
$\rho_t/(\sqrt{2\pi}t)=\sum_{k\in\mathbb Z}\exp(-2\pi^2t^2k^2)$
is decreasing in $t$.

Suppose $\Pr[f\in\mathcal R_q^\times]\ge\tau>0$. For any $0<u<1$,
put $t=u/(8\sigma^2)$. Then the key rejection probability obeys
\begin{equation}\label{eq:tight-key-tail}
	\begin{split}
		\eta_k\le\tau^{-1}\exp\!\left(-tB_k^2+\frac{t}{1-u}\right)
		&\left(\frac{\rho_{\sigma/\sqrt{1-u}}}{\rho_\sigma}\right)^n\\
		{}\cdot&\left(\frac{\rho_{\sigma/\sqrt{1-u/4}}}{\rho_\sigma}\right)^{(D-1)n}.
	\end{split}
\end{equation}
To see this, let $X\leftarrow D_{\mathbb Z,\sigma}$. Completing the
square gives
\[
-\frac{X^2}{2\sigma^2}+t(2X+1)^2
=-\frac{1-u}{2\sigma^2}
\left(X-\frac{u}{2(1-u)}\right)^2+\frac{t}{1-u}.
\]
Only the constant coefficient of $f=2f_0+1$ has the shift by one.
Its exponential moment is therefore at most
$\exp(t/(1-u))\rho_{\sigma/\sqrt{1-u}}/\rho_\sigma$.
Each remaining coefficient of $f$ contributes the same ratio without
that exponential factor; each of the other $(D-1)n$ coefficients
contributes $\rho_{\sigma/\sqrt{1-u/4}}/\rho_\sigma$.
Markov's inequality at $B_k^2$ and division by the invertibility lower
bound $\tau$ prove (1).

For any $0<v<1$, the same exponential-moment argument and the second
theta inequality give
\begin{equation}\label{eq:tight-response-tail}
	1-p_B\le
	\exp\!\left(-\frac{vB_s^2}{2r^2}\right)(1-v)^{-Dn/2}.
\end{equation}
In particular this is an integer-Gaussian bound, rather than substitution
of a continuous chi-square probability.

For the three profiles $q\equiv1\pmod{2n}$ and $\rho_\sigma<2$.
Let $\zeta\in\mathbb F_q$ have order $2n$, let
$H=\langle\zeta\rangle$, write $[x]_q$ for the centered residue of
$x$, and set
\[
E_{n,q}=\min_{a\in\mathbb F_q^\times}
\frac1{2q^2}\sum_{x\in aH}[x]_q^2,
\qquad p_1=\frac{\exp(-1/(2\sigma^2))}{\rho_\sigma}.
\]
Then an admissible invertibility lower bound is
\begin{equation}\label{eq:tight-invertibility}
	\tau=1-\frac nq-n\left(1-\frac1q\right)
	\exp(-16p_1E_{n,q}).
\end{equation}
Indeed the characteristic function
$\phi(s)=\mathbb E\exp(2\pi\mathrm i sX)$ is real, and
$\phi(s)\ge2\Pr[X=0]-1>0$. Retaining the $X=\pm1$ terms gives
\[
0\le\phi(s)\le1-2p_1(1-\cos(2\pi s))
\le\exp(-16p_1\operatorname{dist}(s,\mathbb Z)^2).
\]
For a root $\xi=\zeta^{2j+1}$ of $X^n+1$, finite Fourier inversion
therefore bounds $\Pr[f(\xi)=0]$ by
$q^{-1}+(1-q^{-1})\exp(-16p_1E_{n,q})$.
Here $\{\pm\xi^k:0\le k<n\}=H$ because $n$ is a power of two.
Thus the sum of the $n$ squared distances in each nontrivial Fourier
product is at least $E_{n,q}$; multiplication by two only permutes the
nonzero Fourier frequencies. A union bound over all $n$ roots proves (3).

Exact coset enumeration yields
\[
\begin{array}{c|c|c}
	q&\zeta&E_{n,q}\\ \hline
	12289&10302&11932987670/302039042\\
	15361&4989&18916242006/471920642
\end{array}
\]
The fixed $(u,v)$ values used for profiles 1--3, interpreted as exact
decimal rationals, are
\[
\begin{array}{c|cc}
	&u&v\\ \hline
	1&0.6070989879390658&0.3810220434537782\\
	2&0.5902687398832&0.36167515955237883\\
	3&0.6768302726631384&0.35746618525530693
\end{array}
\]
The outward-rounded supplementary calculation checks
$\eta_k\le2^{-\lambda}$, $1-p_B\le2^{-\lambda}$,
$LM^{-d}(1-2^{-\lambda})\ge\lambda^2$, and
$4U\exp(-\tau(1-2^{-\lambda})L_k)\le2^{-\lambda}$
for the three profiles. This establishes the stated norm and cutoff
inequalities for these distributions. 
	
\end{document}
